\documentclass[10pt]{amsart}
\usepackage[a4paper,margin=23mm]{geometry}
\usepackage{amsmath,amssymb,amsthm,xfrac,hyperref,xspace,graphicx}
\usepackage[english]{babel}
\usepackage[scr=rsfs,cal=euler]{mathalfa}
\newcommand{\resp}{resp.}
\newcommand{\cf}{cf.}
\newcommand{\ie}{i.e.}
\newcommand{\eg}{e.g.}
\newcommand{\loccit}{loc.\ cit.}
\newcommand{\skpt}{\par\smallskip\noindent}
\let\Subsection\subsection
\let\Subsubsection\subsubsection
\emergencystretch=3em
\pagestyle{plain}
\usepackage{aligned-overset}
\hyphenation{defined deri-va-tive ellip-tic encodes every-thing inequal-i-ty oper-a-tor order poten-tial poten-tials solu-tion velo-ci-ty velo-ci-ties}
\usepackage{esint}
\let\strokedint\fint
\multlinegap0pt
\renewcommand{\theenumi}{\alph{enumi}}

\newcommand\mto{\mathchoice{\longmapsto}{\mapsto}{\mapsto}{\mapsto}}
\newdimen\lengtharrow
\newbox\exponentbox

\makeatletter
\def \rightharpoonupfill{$\m@th \mathord-\mkern-6mu \cleaders\hbox{$\mkern-2mu \mathord- \mkern-2mu$}\hfill\mkern-6mu \mathord \rightharpoonup$}
\makeatother

\def\Rightharpoonupts#1%
{\setbox\exponentbox=\hbox{$#1$}
\lengtharrow=\wd\exponentbox
\ifdim \lengtharrow=0pt
		\mathrel{\smash{\mathop{\hbox to 6mm{\rightharpoonupfill}}}}
\else \ifdim\lengtharrow<6mm
\lengtharrow=6mm
\else \advance\lengtharrow by 1mm
\fi
		\overset{\raisebox{3pt}{$\scriptstyle#1$}}{\mathrel{\smash{\hbox to\lengtharrow{\rightharpoonupfill}}}}
	\fi}
	
\def\Rightharpoonupds#1%
{\setbox\exponentbox=\hbox{$#1$}
\lengtharrow=\wd\exponentbox
\ifdim \lengtharrow=0pt
		\mathrel{\smash{\mathop{\hbox to 6mm{\rightharpoonupfill}}}}
\else \ifdim\lengtharrow<6mm
\lengtharrow=6mm
\else \advance\lengtharrow by 1mm
\fi
		\overset{\raisebox{5pt}{$\scriptstyle#1$}}{\mathrel{\smash{\hbox to\lengtharrow{\rightharpoonupfill}}}}
	\fi}
	
\def\Rightharpoonup#1{\mathchoice{\Rightharpoonupds{#1}}{\Rightharpoonupts{#1}}{\Rightharpoonupts{#1}}{\Rightharpoonupts{#1}}}

\let\epsilon\varepsilon
\let\backslash\setminus
\let\ts\textstyle
\newcommand{\Psfrac}[2]{(\sfrac{#1}{#2})}
\newcommand{\psfrac}[2]{\sfrac{(#1)}{#2}}
\newcommand{\spfrac}[2]{\sfrac{#1}{(#2)}}
\newcommand{\pspfrac}[2]{\sfrac{(#1)}{(#2)}}

\newcommand{\const}{\textup{const}}
\newcommand{\constants}{\textup{constants}}
\newcommand{\loc}{\mathrm{loc}}
\newcommand{\sbullet}{{\scriptscriptstyle\bullet}}
\newcommand{\rI}{\mathrm{I}}
\newcommand{\rII}{\mathrm{II}}
\newcommand{\rIII}{\mathrm{III}}
\newcommand{\rIV}{\mathrm{IV}}
\newcommand{\rV}{\mathrm{V}}
\newcommand{\rVI}{\mathrm{VI}}


\newtheorem{theorem}{Theorem}[subsection]
\newtheorem{corollary}[theorem]{Corollary}
\newtheorem{proposition}[theorem]{Proposition}
\newtheorem{lemma}[theorem]{Lemma}
\newtheorem{condition}[theorem]{Condition}

\theoremstyle{definition}
\newtheorem{definition}[theorem]{Definition}
\newtheorem{remark}[theorem]{Remark}
\newtheorem{example}{Example}[subsection]

\renewcommand\star{\textsuperscript{*} }
\renewcommand{\Re}{\operatorname{Re}}
\renewcommand{\Im}{\operatorname{Im}}
\renewcommand{\div}{\operatorname{div}}
\DeclareMathOperator{\dist}{dist}
\DeclareMathOperator{\curl}{curl}

\newcommand\R{\mathbb{R}}
\newcommand\B{\mathbf{B}}
\newcommand\T{\mathbb{T}}
\newcommand\Z{\mathbb{Z}}
\newcommand\sym{\mathcal{E}}
\newcommand\N{\mathbb{N}}

\newcommand\norm[1]{\left\lVert #1\right\rVert}
\newcommand\anorm[1]{\lVert #1\rVert}
\newcommand\bnorm[1]{\bigl\lVert #1\bigr\rVert}
\newcommand\Bnorm[1]{\Bigl\lVert #1\Bigr\rVert}
\newcommand\bbnorm[1]{\biggl\lVert #1\biggr\rVert}
\newcommand\scalar[2]{\langle #1,#2\rangle}
\newcommand\bscalar[2]{\bigl\langle #1,#2\bigr\rangle}
\newcommand\Bscalar[2]{\Bigl\langle #1,#2\Bigr\rangle}

\newcommand\intT{\int_{\T^d}}
\newcommand\intR{\int_{\R^d}}
\newcommand\eq{=}

\newcommand\dd{\,\mathrm{d}}
\newcommand\dx{\dd x}
\newcommand\dz{\dd z}
\newcommand\dR{\dd R}
\newcommand\dZ{\dd Z}
\newcommand\dy{\dd y}
\newcommand\dt{\dd t}
\newcommand\ds{\dd s}
\newcommand\dr{\dd r}

\newcommand{\mres}{\mathbin{\vrule height 1.6ex depth 0pt width
0.13ex\vrule height 0.13ex depth 0pt width 1.3ex}}

\DeclareFontFamily{U}{mathx}{}
\DeclareFontShape{U}{mathx}{m}{n}{<-> mathx10}{}
\DeclareSymbolFont{mathx}{U}{mathx}{m}{n}
\DeclareMathAccent{\widehat}{0}{mathx}{"70}
\DeclareMathAccent{\widecheck}{0}{mathx}{"71}

\let\chek\widecheck
\let\de\partial
\let\eps\epsilon
\let\mel\relax

\begin{document}
\title{Finite-dimensional added-inertia ODE reduction for volume-preserving toroidal solids in axisymmetric irrotational no-swirl Euler flow}
\author{David Meyer}
\maketitle
\section{The source fluid-solid model}
Mathematically, the setup for this is as follows: We use axisymmetric coordinates $(r,z)$ and denote the right half-plane by $\mathbb{H}$. We~fix the number of filaments $k\in \N_{>0}$ and numbers $v_1,\dots, v_k>0$, which we interpret as the volumes of the bodies and which should remain constant along the evolution. We~set
\begin{align}
\rho_i:=\sqrt{\frac{v_i}{\pi R_i}},\qquad B_i(R_i,Z_i):=B_{\rho_i}((R_i,Z_i))\subset\mathbb{H} \quad\text{for $i\eq1,\dots,k$,}\label{def rho}
\end{align}
here $\rho_i$ denotes the minor radius, $R_i$ the major radius and $Z_i$ the $Z$-coordinate.
Formally, we~describe the configuration of the bodies through the manifold $M\subset (\R^2)^k$ of all $(R_1,Z_1),\dots, (R_k,Z_k)$ such that the bodies $B_i$ all have positive distance from each other and from $\de\mathbb{H}$.

For $q\in M$, we~shall write $q_i$ for the $i$-th component and $q_{R_i}$ and $q_{Z_i}$ for the two components of $q_i$. We~shall also write $B_i(q)$ for clarity instead of $B_i$ sometimes. Let~$n$ denote the \textit{outer} normal of $\bigcup_i\de B_i$.

Each body can undergo two different kinds of motion, the rigid one in the $Z$\nobreakdash-direction and the non-rigid change of length. To make this precise, we~use the natural correspondence between tangent vectors and normal velocities on the $\de B_i$, as~every $C^1$-curve in $M$ corresponds to a continuous movement of each $B_i$.

For a tangent vector $t^*$, let $t_{R_1}^*,t_{Z_1}^*,t_{R_2}^*,\dots$ denote its components. We~say that a tangent vector is associated with $B_i$ if only its $R_i$- and $Z_i$-component are non-zero, and write $T_{q_i}M$ for the subspace of those tangent vectors.

Then, for a $C^1$-curve $q$, the normal velocity is given by
\begin{align}
u(\dot{q}):=\dot{q}_{Z_i}n\cdot e_Z+\dot{q}_{R_i}n\cdot e_R-\frac{\dot{q}_{R_i}\rho_i}{2R_i}\quad\text{on $\de B_i$.}\label{normal velo}
\end{align}
Here $e_R$ and $e_Z$ are the unit vectors and the purpose of the last summand is to make sure that if the major radius changes, the minor radius also changes so that the volume of $B_i$ is conserved under the motion.

Indeed by using Gauss's theorem, we~see
\begin{equation}
\begin{aligned}
\int_{\de B_i}ru(\dot{q})\dx&=\int_B \div\left(r(\dot{q}_{Z_i}e_Z+\dot{q}_{R_i}e_R)\right)\dx-\frac{\dot{q}_{R_i}\rho_i}{2R_i}\int_{\de B_i} r\dx\\
&=\pi\rho_i^2\dot{q}_{R_i}-\pi \rho_i^2\dot{q}_{R_i}=0,\label{vol fix}
\end{aligned}
\end{equation}
and hence by the Reynolds transport theorem this normal velocity keeps the volume (with respect to $r\dr\dz$) fixed.

Let
\begin{align*}
\mathcal{F}=\mathcal{F}(t):=\mathbb{H}\backslash \bigcup_i B_i
\end{align*}
denote the (time-dependent) domain of the fluid. Let $L_R^2$ and $H_R^1$ denote the $L^2$ \resp $H^1$ spaces with respect to the measure $r\dr\dz$. We~assume that:

\begin{condition}
In $\mathcal{F}(t)$ the fluid fulfills the axisymmetric Euler equations with zero vorticity and no swirl for $t\in [0,\infty)$, \ie
\begin{align}
\de_t u+(u\cdot\nabla)u+\nabla p&\eq 0,\label{Euler1}\\
\div(ru)&\eq 0,\label{Euler2}\\
\curl(u)&\eq 0,\label{Euler3}\\
u_r&\eq 0 \quad\text{\normalfont{on} $\de \mathbb{H}$.}\label{Euler4}
\end{align}
Here the $\div$ and $\curl$ are taken with respect to the variables $(r,z)$ and $u$ is $\R^2$-valued.
\end{condition}

These equations are equivalent to the usual Euler equations for $u$ with no vorticity and no swirl (that is, no velocity in the direction perpendicular to the $(r,z)$-plane) after going back to three-dimensional coordinates. The boundary condition~\eqref{Euler4} encodes the fact that there should be no singularity at $\{r=0\}$.

\begin{condition}\label{strongsol}
We assume that the solution of \eqref{Euler1}--\eqref{Euler4} is strong in the sense that $u,\nabla u,\de_t u\in L_R^2\cap C^1(\mathcal{F})$ and that $q$ is $C^2$ in time. In~particular this should hold for the initial data.
\end{condition}

We define the circulation around each body as
\begin{align*}
\gamma_i:=\int_{\de B_i}u\cdot\tau\dx,
\end{align*}
where $\tau=n^\perp$. This is a conserved quantity by Kelvin's circulation law. For technical reasons, we~will assume:
\begin{condition}\label{cond:103}
None of the $\gamma_i$ are $0$.
\end{condition}

This condition is not necessary for the well-posedness (nothing in the proof changes),
but necessary for the convergence to the limit system, for instance it follows directly from the
ODE reformulation \eqref{MainODE} that a single body with no circulation which is initially moving in the $e_z$-direction will keep moving in that direction with the same speed as every term
in \eqref{MainODE} is 0 in that case.

We assume that the normal velocity of the boundary of each $B_i$ matches the velocity of $u$ in the corresponding direction:
\begin{condition} For all $i$ we have
\begin{align}
u\cdot n\eq u(\dot{q}) \quad\text{on $\de B_i$}\label{Euler5},
\end{align}
where we use the identification between tangent vectors and normal velocities mentioned above.
\end{condition}

For simplicity we will assume that all bodies and the fluid have constant density~$1$, though all the arguments still work for different densities.

\begin{condition} We assume that in the $z$-direction, the momentum is (formally) preserved, which yields the condition
\begin{align}
v_i\ddot{q}_{{Z}_i}\eq-\int_{\de B_i}rpn\cdot e_Z\dx\label{eq z}
\end{align}
for all $i$.
\end{condition}

It remains to derive a condition for the interaction of the fluid and the solids in the $r$-component. As the solids can change their shape, we~make the Ansatz of prescribing an interior velocity field and assuming that the kinetic energy of each $B_i$ only changes through the force exerted by the pressure at the boundary.

We associate to each tangent vector/normal velocity $t_i^*$ associated with $B_i$ such an interior velocity field $u_{i,\mathrm{int}}(t_i^*)\in H^1(B_i)$ with
\begin{align}
\div ru_{i,\mathrm{int}}&\eq0 \quad\text{in $B_i$},\label{Int1}\\
\curl u_{i,\mathrm{int}}&\eq0 \quad\text{in $B_i$},\label{Int2}\\
u_{i,\mathrm{int}}\cdot n&\eq u(t_i^*)\quad\text{on $\de B_i$.}\label{Int3}
\end{align}
Existence and uniqueness can be obtained by standard elliptic theory \cite[Chap.\,6]{evans2022partial}, as~$u_{i,\mathrm{int}}(t_i^*)$ can be written as $\nabla \phi$ with $\div(r\nabla \phi)=0$ and the boundary condition $\de_n \phi=u(t_i^*)$ by the assumption of curl-freeness. The compatibility condition for the Neumann boundary condition $\int_{\de B_i} r\de_n \phi\dx=0$ is fulfilled since $\int_{\de B_i} r u(t_i^*)\dx=0$ by \eqref{vol fix}.
Note that this is linear in $t_i^*$, and that
\begin{align}\label{Ztrivial}u_{i,\mathrm{int}}(t_{Z_i}^*)=t_{Z_i}^*e_Z.
\end{align}
We can use this to associate a quadratic form on $T_{q_i}M$ by
\begin{align}
(t_i^*)^TE_{q_i}t_i^+:=\int_{B_i}r\scalar{u_{i,\mathrm{int}}(t_i^*)}{u_{i,\mathrm{int}}(t_i^+)}\dx, \label{def E}
\end{align}
for tangent vectors $t_i^*,t_i^+$, associated with $B_i$, which describes the kinetic energy associated with $u_{i,\mathrm{int}}$. Clearly, this is symmetric and positive definite. Because of the explicit form \eqref{Ztrivial}, we~know that
\begin{align*}
(e_Z)_i^TE_{q_i}(e_Z)_i\eq v_i,
\end{align*}
where $(e_Z)_i\in (\R^2)^k$ is the vector with $e_Z$ in the $i$-th component. We~set
\begin{align}
(e_R)_i^TE_{q_i}(e_R)_i=: f_i(q_{R_i})\label{def f},
\end{align}
here the function $f_i$ depends on the volume $v_i$ and $(e_R)_i$ is the vector with $e_R$ in the $i$-th component.

In summary, the quadratic form can be written in components as
\begin{align}
(t_i^*)^T\begin{pmatrix} f_i(q_{R_i}) & 0\\ 0 & v_i\end{pmatrix}(t_i^+),\label{def f 2}
\end{align}
as one can see from the explicit formula \eqref{Ztrivial} and the fact that $u_{i,\mathrm{int}}(t_{R_i}^*)$ must be asymmetric in $z$ in the second component and hence $t_{R_i}^*$ and $t_{Z_i}^*$ are orthogonal to each other with respect to $E_{q_i}$.

We define the kinetic energy of $B_i$ as
\begin{align}\label{intEnergy}
\mathcal{E}_{B_i}:=\frac{1}{2}\dot{q}_i^TE_{q_i}\dot{q}_i.
\end{align}
We assume that the kinetic energy of each $B_i$ only changes through the force exerted at the boundary, that is,
\begin{align}
\mathcal{E}_{B_i}'\eq-\int_{\de B_i} rpu(\dot{q}_i)\dx.\label{energyBalance}
\end{align}
After using the decomposition \eqref{def f 2} and subtracting the condition \eqref{eq z}, one obtains that
\begin{align*}%\label{eq r}
f_i(q_{R_i})\ddot{q}_{R_i}\dot{q}_{R_i}+\frac{1}{2}\de_{q_{R_i}}f_i(q_{R_i})(\dot{q}_{R_i})^3\eq -\int_{\de B_i}r p\dot{q}_{R_i}\Bigl(n\cdot e_R -\frac{\rho_i}{2R_i}\Bigr)\dx.
\end{align*}
We make the extra assumption that one can divide out $\dot{q}_{R_i}$, which is equivalent to saying that whenever $\dot{q}_{R_i}=0$ and the force at the boundary is nonzero, then $\ddot{q}_{R_i}$ is not zero, which rules out bodies with fixed $R_i$-coordinate. This gives the final equation:
\begin{condition}
For all $i$ we have
\begin{align}\label{eq R}
f_i(q_{R_i})\ddot{q}_{R_i}+\frac{1}{2}\de_{q_{R_i}}f_i(q_{R_i})(\dot{q}_{R_i})^2\eq -\int_{\de B_i}r p \Bigl(n\cdot e_R -\frac{\rho_i}{2R_i} \Bigr)\dx.
\end{align}
\end{condition}

We can also write \eqref{eq z} and \eqref{eq R} as a single equation
\begin{align}\label{2ndDeri}
(t_i^*)^TE_{q_i}\ddot{q}_i+\frac{1}{2}t_i^*(\de_{\dot{q}_i} E_{q_i}\cdot \dot{q}_i)\dot{q}_i\eq-\int_{\de B_i} rpu(t_i^*)\dx,
\end{align}
where $t_i^*$ is an arbitrary tangent vector associated with $B_i$.
If $\Omega\subset \mathbb{H}$ we write $L_0^2(\Omega)$ for the space of all functions $f\in L^2(\Omega)$ such that for every connected component $\Omega_i$ of $\Omega$ we have $\int_{\Omega_i}f\dx=0$.

If $S\subset \mathbb{H}$, we~write $S^{\R^3}$ for its figure of revolution in $\R^3$.
\section{Well-posedness of the system}\label{Section2}
We will follow the approach in \cite{GlassMunnierSueur2} to show that our system can be reduced to an ODE in $q$ only (see Theorem \ref{ODEtheorem} below), which in particular shows well-posedness. The main additional difficulty is that we are in an unbounded domain and hence need decay estimates to justify partial integrations.

\subsection{Representation of $u$}\label{repU}
In this subsection, we~show that $u$ can be recovered from a potential and a streamfunction (defined in \ref{def psi}), which is shown in Proposition~\ref{prop rep u} below.

\begin{definition}
For $t^*\in T_{q_i}M$, let $\phi_{i,t^*}=\phi_{i,t^*}(t)$ be defined as the unique solution of the Neumann problem
\begin{align*}
\div(r\nabla\phi_{i,t^*})&\eq0 \quad\text{\normalfont{in} $\mathcal{F}(t)$},\\
\de_n\phi_{i,t^*}&\eq u(t^*) \quad\text{\normalfont{on} $\de B_i$},\\
\de_n\phi_{i,t^*}&\eq 0 \quad\text{\normalfont{on} $\de B_j$ \normalfont{for} $j\neq i$ \normalfont{and on} $\de\mathbb{H}$},\\
\phi_{i,t^*}&\in \dot{H}_R^1,\\
\phi_{i,t^*}&\to 0\quad\text{at $\infty$}.
\end{align*}
\end{definition}

We first need to check that this is well-defined.

\begin{lemma}\label{ExistencePhi}
Let $b\in L^2(\cup_i \de B_i)$ be such that $\int_{\cup_i\de B_i}rb\dx=0$, then the equations
\begin{align*}
\div(r\nabla\phi)&\eq0 \quad\text{\normalfont{in} $\mathcal{F}(t)$},\\
\de_n\phi_{}&\eq b \quad\text{\normalfont{on} $\ts\bigcup_i \de B_i$},\\
\de_n\phi_{}&\eq 0 \quad\text{\normalfont{on} $\de\mathbb{H}$},\\
\phi_{}&\in \dot{H}_R^1,\\
\phi_{}&\to 0\quad\text{at $\infty$}
\end{align*}
have a unique solution $\phi$. Furthermore
\begin{align}\label{decayphi}
|\nabla^m\phi(x)|\lesssim \frac{\norm{b}_{H^{m}}}{1+|x|^{2+m}} \quad\forall m\in\N_{\geq 0},
\end{align}
where the implicit constant is bounded locally uniformly in $q$.

\end{lemma}

This implies the well-definedness of $\phi_{i,t^*}$ and that the estimate \eqref{decayphi} holds locally uniformly in $q$ for $\phi_{i,t^*}$, since $\int_{\cup_i\de B_i} ru(t^*)\dx=0$, as~computed in \eqref{vol fix}.

\begin{proof}
We go back to three-dimensional coordinates and set $\phi^{\R^3}(r,z,0)=\phi(r,z)$, where $\phi^{\R^3}$ is axisymmetric. Then $\phi$ solves the system above iff $\phi^{\R^3}$ solves the corresponding Neumann problem for $\Delta$ in $\R^3\backslash \bigcup_j B_j^{\R^3}$. By~standard techniques (see \eg \cite{Amrouche}), we~obtain a unique solution $\phi^{\R^3}\in \dot{H}^{1}(\R^3\backslash \bigcup B_j^{\R^3})$. We~furthermore obtain from this that $\de_n\phi=0$ on $\de\mathbb{H}$.

The decay rate then directly follows from the lemma below.
\end{proof}

\skpt
\begin{lemma}\label{decay1}
\begin{enumerate}
\item\label{decay1a}
Let $\zeta\in \dot{H}^1(\mathcal{F}^{\R^3})$ be axisymmetric such that $\Delta\zeta=0$ in $\mathcal{F}^{\R^3}$ and $\zeta(x)\to 0$ as $|x|\to\infty$. Then it holds that
\begin{align*}
\left|\nabla^m\zeta(x)\right|\lesssim \frac{\norm{\de_n\zeta}_{H^{m}(\de\mathcal{F}^{\R^3})}}{1+|x|^{1+|m|}}\qquad \forall m\in\N_{\geq 0}.
\end{align*}
The implicit constant is bounded locally uniformly in $q$.

\item\label{decay1b}
 Let $\zeta$ be as in \eqref{decay1a}. If~additionally
\begin{align*}
\int_{\cup_j \de B_j^{\R^3}}\de_n\zeta\dx\eq0,
\end{align*}
then
\begin{align*}
\left|\nabla^m\zeta(x)\right|\lesssim \frac{\norm{\de_n\zeta}_{H^{m}(\de\mathcal{F}^{\R^3})}}{1+|x|^{2+|m|}}\qquad \forall m\in\N_{\geq 0},
\end{align*}
where the implicit constant is controlled as in \eqref{decay1a}.
\end{enumerate}
\end{lemma}

\begin{proof}
If we extend $\zeta$ to $\R^3$ by solving the Dirichlet problem for $\zeta|_{\de B_j^{\R^3}}$ on each~$B_j^{\R^3}$, then the (distributional) Laplacian of this extension is of the form
\[
[\de_n\zeta]\mathcal{H}^2\mres{ \bigcup \de B_j^{\R^3}},
\]
 where $[\cdot]$ denotes the jump across the boundary, as~a direct calculation shows. By~elliptic regularity (\cf \cite[Chap.\,2]{grisvard2011elliptic}), this is a finite measure.

We claim that we can recover this extension by convoluting the distributional Laplacian with the Newtonian potential. Indeed for any $f\in C_c^\infty(\R^3)$, we~have
\begin{align*}
\int_{\R^3}\int_{\cup_i\de B_i^{\R^3}}\frac{-[\de_n\zeta]}{4\pi|x-y|}\Delta f(x)\,\text{d}\mathcal{H}^2(y)\dx\eq\int_{\cup_i\de B_i^{\R^3}}[\de_n\zeta]f(y)\,\text{d}\mathcal{H}^2(y)\eq \int_{\R^3}\zeta \Delta f\dx,
\end{align*}
where in the second step, we~used that the Newtonian potential is the inverse Laplacian. Hence the difference $[\de_n\zeta]\mathcal{H}^2\mres(\bigcup_i\de B_i^{\R^3})*\Psfrac{-1}{4\pi|x|}-\zeta$ is a harmonic tempered distribution, \ie a polynomial.

Now
\begin{align*}
[\de_n\zeta]\mathcal{H}^2\mres({\ts\bigcup_i\de B_i^{\R^3}})*\frac{-1}{4\pi|x|}-\zeta\to 0,
\end{align*}
by the assumption on $\zeta$ and because $[\de_n\zeta]\mathcal{H}^2\mres(\bigcup_i\de B_i^{\R^3})$ is a finite measure, so this difference is zero, which shows the claim.

We hence obtain that $|\nabla^m\zeta(x)|\lesssim\spfrac{\norm{\de_n\zeta}_{L^2(\de\mathcal{F}^{\R^3})}}{1+|x|^{1+m}}$ for all $m\geq0$ and for $\dist(x,\bigcup B_i^{\R^3})\geq 1$. For $x$ close to the $B_j$ the estimate follows from elliptic regularity theory. This shows the estimate in part \eqref{decay1a}, part \eqref{decay1b} works exactly the same way, except that the integral of $[\de_n\zeta]$ vanishes by the assumption and partial integration, which gives one order more of decay.

To see that these bounds are locally uniform in $q$, we~need uniform estimates on $\norm{[\de_n\zeta]}_{L^1(\cup_i\de B_i)}$. Note that for this we only need up-to-the-boundary estimates for a neighborhood of each $B_i$, locally uniform in $q$ and a locally uniform $L^2$ estimate. By~going back to axisymmetric coordinates, one obtains the former, as~the geometries of these neighborhoods (in axisymmetric coordinates) only change by rescaling with a bounded factor.
On the other hand, one can obtain from the energy equality (which is justified by the decay estimates we have already proved) in axisymmetric coordinates that
\begin{align}\label{energy eq}
\int_{\mathcal{F}} r|\nabla\zeta|^2\dx\eq -\int_{\cup_i \de B_i}r\zeta\de_n \zeta\dx\gtrsim \anorm{\zeta}_{H^1(\cup_i (B_i+B_1(0)\backslash B_i))}^2.
\end{align}
Here we used the (three-dimensional) Sobolev inequality to control the $L^2$-norm on the right-hand side. The constant of the Sobolev inequality is locally uniform in $q$, as~one can \eg see by using a diffeomorphism between the different instances of $\mathcal{F}$.

Now we can use the trace inequality from $H^1(B_i+B_1(0)\backslash B_i)$ to $L^2(\de B_i)$ in \eqref{energy eq}, whose constant is also locally uniform in $q$, as~the geometry only changes through rescaling by a bounded factor. This implies the desired locally uniform estimate
\[
\norm{\de_n\zeta}_{L^2(\cup_i\de B_i)}\gtrsim\anorm{\zeta}_{L^2(\cup_i (B_i+B_1(0)\backslash B_i))}.\qedhere
\]
\end{proof}

\begingroup

The argument for the existence of the stream function is more complicated, as~there is no easy algebraic relation between the three-dimensional stream function and the axisymmetric stream function.

\begin{definition}\label{def psi}
Let $\psi_{i}=\psi_i(t)$ be the solution of the elliptic equation
\begin{equation}
\label{Flow}
\begin{aligned}
\div\Bigl(\frac{1}{r}\nabla\psi_{i}\Bigr)&\eq0\quad \mathrm{ in }\quad\text{}\mathcal{F}(t),\\[-3pt]
\psi_{i}|_{\de B_j}&\text{ is constant $\forall j$},\\[-3pt]
\int_{\de B_j}\frac{1}{r}\de_{n}\psi_{i}\dx&\eq\delta_{ij},\\[-3pt]
\psi_i|_{\{r=0\}}&\eq0,\\[-3pt]
\lim_{r\to 0}\frac{1}{r}\de_z\psi_i&\eq 0.
\end{aligned}
\end{equation}
Here $\delta_{ij}$ is a Kronecker delta. We~refer to the constant boundary values on $\de B_j$ as~$C_{ij}$.
\end{definition}

\endgroup

\begin{lemma}\label{ExistencePsi}
Such a $\psi_i$ always exists and is unique under the constraint $\frac{1}{\sqrt{r}}\nabla\psi_i\in\nobreak L^2$. Furthermore $\frac{1}{r}\nabla\psi_i$ is continuous at $r=0$ and we have the estimate\vspace*{-3pt}
\begin{align*}
\Bigl|\nabla^m\Bigl(\frac{1}{r}\nabla\psi_i(x)\Bigr)\Bigl|\lesssim\frac{1}{1+|x|^{2+m}}\quad \forall m\in\N_{\geq 0}.
\end{align*}
The implicit constant in the estimate is locally uniformly bounded in $q$.
\end{lemma}
\begin{proof}
We first show that an auxiliary function $u_{2,i}$ can be constructed by going back to three-dimensional coordinates and later show that one can recover $\psi_i$ from it.

\subsubsection*{Step 1}
We define $u_{2,i}$ on $\mathcal{F}^{\R^3}$ as the axisymmetric vector field with no azimuthal component which solves (in three-dimensional variables) the system
\begin{align}
\div u_{2,i}&\eq0,\label{DivCurl}\\[-3pt]
 \curl u_{2,i}&\eq0,\label{DivCurl3}\\[-3pt]
 u_{2,i}\cdot n&\eq0\label{DivCurl4}\quad\text{on $\de B_j^{\R^3}$ for all $j$},\\[-3pt]
 u_{2,i}&\in L^2(\mathcal{F}^{\R^3}).\label{DivCurl5}
\end{align}
We make the ansatz $\curl \Psi_i=u_{2,i}$ for a purely azimuthal field $\Psi_i=\tilde{\Psi}_ie_\theta$, where $e_\theta$ is the unit vector in the $\theta$-direction.
This gives the equations
\begin{align*}
(u_{2,i})_r\eq-\de_z\tilde{\Psi}_i,\quad (u_{2,i})_z\eq\de_r\tilde{\Psi}_i+\frac{1}{r}\de_r \tilde{\Psi}_i,
\end{align*}
using that the last term can be rewritten as $\frac{1}{r}\de_r(r\Psi_i)$, this turns the equations \eqref{DivCurl}--\eqref{DivCurl4} into the equations
\begin{gather}
\Delta \Psi_i\eq0\quad\text{in $\mathcal{F}^{\R^3}$},\label{BigPhi1}\\[-3pt]
r\tilde{\Psi}_i|_{\de B_j^{\R^3}}\text{ is constant for all $j$.}\label{BigPhi2}
\end{gather}
For each fixed set of constant boundary values $(\tilde{C}_{ij})_{j=1,\dots,k}$ for $r\tilde{\Psi}_i$, this system has a unique solution $\Psi_i(\tilde{C}_{ij})\in H^1$ by standard techniques. Also, because $\Psi_i$ is purely azimuthal, it must vanish at $r=0$ and hence it holds that\vspace*{-3pt}
\begin{align}\label{u at 0}
(u_{2,i})_r\eq0
\end{align}
at $r=0$.

\subsubsection*{Step 2} To show uniqueness of $u_{2,i}$ for given $(\tilde{C}_{ij})$, we~note that we can recover a~$\Psi_i$ fulling the system \eqref{BigPhi1},\eqref{BigPhi2} from $u_{2,i}$. Indeed we may extend $u_{2,i}$ to the full space by zero as $\overline{u}_{2,i}$, which preserves the condition that the divergence vanishes, as~the distributional divergence on the boundary equals $[\overline{u}_{2,i}\cdot n]=0$.

It is well known that on the full space, every divergence-free field can be written as a curl, indeed we have $-\curl\Delta^{-1}\curl g=g$ for every divergence-free $g$, as~a straightforward calculation shows.
Since we also have $\curl \overline{u}_{2,i}\in H^{-1}$ and the fundamental solution of the Laplacian maps $H^{-1}$ to $H^{1}$, we~see that the field obtained this way lies in $H^1$.
Hence two different solutions $u_{2,i}$ for the same boundary values would give rise to two different $\Psi_i$ in $H^1$ with the same boundary values, which is impossible.

\subsubsection*{Step 3} Now we can view $u_{2,i}$ as a function in $(r,z)$ again, it fulfills $\div(ru_{2,i})=0$ (with the two-dimensional divergence). Then we can find a $\psi_i(\tilde{C}_{ij})$ such that
\begin{align*}
u_{2,i}\eq\frac{1}{r}\nabla^\perp\psi_i,
\end{align*}
because $ru_{2,i}^\perp$ is curl-free. This can be done with the usual path integral construction, it is easy to check that the condition $u_{2,i}\cdot n=0$ ensures that even paths which are not homotopy equivalent yield the same values. One can then check by direct calculation that $\div(\frac{1}{r}\nabla\psi_i(\tilde{C}_{ij}))=0$ holds, and by the boundary condition for $u_{2,i}$, we~see that $\psi_i(\tilde{C}_{ij})$ must be constant on each $\de B_j$. Furthermore, we~have that $u_{2,i}$ is continuous at $r=0$ by elliptic regularity and hence $\frac{1}{r}\nabla\psi_i(\tilde{C}_{ij})$ is continuous at $r=0$ and $\frac{1}{r}\de_z\psi_i=0$ at $r=0$ by \eqref{u at 0}.
This $\psi_i(\tilde{C}_{ij})$ is unique up to an additive constant under the condition $\frac{1}{\sqrt{r}}\nabla\psi_i(\tilde{C}_{ij})\in L^2$ (here one gets an additional factor $r$ from the coordinate change), as~one can recover the unique $u_{2,i}$ from it.

Next, we~argue that we can pick the boundary values $(\tilde{C}_{ij})$ uniquely such that the condition \eqref{Flow} holds. It suffices to show that the linear map that sends the boundary values $(\tilde{C}_{ij})$ to the integrals $\int_{\de B_j}\frac{1}{r}\de_{n}\psi_i(\tilde{C}_{ij})\dx$ is invertible for each fixed $i$.

First note that the $\tilde{C}_{ij}$ are also the boundary values of $\psi_i$ (up to an additive constant) because we have that $ru_{2,i}=\nabla^\perp r\tilde{\Psi}_i$ (in $(r,z)$-coordinates) and hence by applying the fundamental theorem of calculus, we~see that $(r\tilde{\Psi})(x)-(r\tilde{\Psi})(y)=\psi_i(x)-\psi_i(y)$ (in axisymmetric coordinates).

Assume there is a nonzero vector of $\tilde{C}_{ij}$ such that all integrals vanish. Without loss of generality, we~may assume that $\tilde{C}_{i1}$ is the biggest one of the $\tilde{C}_{ij}$. Then the normal derivative of $\psi_i$ on $\de B_1$ must be non-positive by the maximum principle and hence must be zero everywhere on $\de B_1$. Since the tangential derivative also vanishes, the constant extension of $\psi_i(\tilde{C}_{ij})$ to $B_1$ still fulfills $\div(\frac{1}{r}(\psi_i(\tilde{C}_{ij}))=0$. But this extension is then a locally, but not globally, constant solution of an elliptic equation, which is a contradiction.

We have $\de_z\psi_i=0\cdot u_{2,i}=0$ at $r=0$ because $u_{2,i}$ is continuous by elliptic regularity. We~may choose the additive factor that we have leftover such that $\psi_i=0$ at $r=0$ holds.

The uniqueness of $\psi_i$ follows from the uniqueness of the $\psi_i(\tilde{C}_{ij})$.

\subsubsection*{Step 4} By elliptic regularity, it is easy to see that the $C_{ij}$ are locally uniformly bounded in $q$, and hence $\norm{\de_n\Psi_i(C_{ij})}_{H^m}$ is locally uniformly bounded in $q$. By~Lem\-ma~\ref{decay1}\,\eqref{decay1a}, we~see that
\begin{align*}|\nabla^m\Psi_i(x)|\lesssim \frac{1}{1+|x|^{1+m}},
\end{align*}
for all $m\in \N_{\geq 0}$ which implies the decay estimate.
\end{proof}

It is known that the equations
\begin{align*}
\div\Bigl(\frac{1}{r}\nabla f\Bigr)\eq g\quad\text{in $\mathbb{H}$},\\
f\eq\de_nf\eq 0\quad\text{on $\de\mathbb{H}$},
\end{align*}
have a fundamental solution $K$ such that
\begin{align}\label{definition K}
f(y)\eq \int_{\mathbb{H}}g(x)K(x,y)\dx
\end{align}
is the unique solution under suitable decay assumptions on $f$ for \eg $g\in C_c^\infty(\mathbb{H})$, see \eg \cite[\S 2]{GallaySverak}.

By the following lemma, we~can recover solutions to $\div(\frac{1}{r}\nabla\cdot)$ from single-layer potentials with this fundamental solution and hence obtain decay estimates.

\skpt
\begin{lemma}\label{rep psi}
\begin{enumerate}
\item\label{rep psia}
Let $\div(\frac{1}{r}\nabla\zeta)=0$ in $\mathcal{F}$ with $\frac{1}{\sqrt{r}}\nabla\zeta\in L^2(\mathcal{F})$, assume that $\frac{1}{r}\nabla\zeta$ is continuous for $r\to 0$ and that $\zeta|_{r=0}=0$. Furthermore, assume that $\zeta|_{\de\mathcal{F}}$ is sufficiently smooth, then there is a constant $C$, depending on $\zeta$ such that
\begin{align*}
|\zeta(x)|\leq \frac{C}{1+|x|}.
\end{align*}
In particular this holds for $\psi_i$.
\item\label{rep psib}
$\psi_i$ can be represented as
\begin{align*}
\psi_i(y)\eq\int_{\cup_i\de B_i}\frac{1}{r}K(x,y)\de_n\psi_i\dx.
\end{align*}
\end{enumerate}
\end{lemma}
\begin{proof}
%\eqref{rep psia} and \eqref{rep psib}
We claim that if we extend $\zeta$ to $\mathbb{H}$ by solving the Dirichlet problem for $\div(\frac{1}{r}\nabla\cdot)$ with boundary values $\zeta$ in each $B_i$, then it holds that
\begin{align*}
\zeta\eq \int_{\cup_i\de B_i}\frac{1}{r}K(x,y)[\de_n\zeta]\dx.
\end{align*}
This directly shows \eqref{rep psib} because for $\psi_i$ the extension to each $B_i$ is constant. It also shows \eqref{rep psia} by the fact that $[\de_n\zeta]\mathcal{H}^1\mres \de B_i$ is a finite measure by elliptic regularity and the fact that the fundamental solution $K(x,y)$ decays like $\spfrac{1}{1+|y|}$ at $\infty$ locally uniformly in $x$ (see \cite[Lem.\,2.1\,ii)]{GallaySverak}).

The same argument as in the proof of Lemma \ref{decay1} shows that
\begin{align*}
g(y):=\int_{\de\mathcal{F}}\frac{1}{r}K(x,y)[\de_n\zeta]\dx
\end{align*}
fulfills $\div(\frac{1}{r}\nabla (g-\zeta))=0$. Furthermore, by the aforementioned decay estimates for~$K$ we see that $|g(x)|\lesssim \spfrac{1}{|1+|x|}$, that $g=0$ at $r=0$ and that $\frac{1}{r}\nabla g$ is continuous at~$0$. By~elliptic regularity, it holds that $g-\zeta$ is smooth in the interior of $\mathbb{H}$.

Now let $h(r,z,\theta)=(\frac{1}{r}\nabla(\zeta-g)(r,z))^\perp$ (in axisymmetric coordinates). This function is divergence-free (with respect to three-dimensional variables) as a direct calculation shows for $r>0$, at $r=0$ it is also divergence-free by continuity.
Hence there is an~$H$ with $\curl H=h$, where the gradient is taken with respect to three-dimensional variables. Then it holds that $\Delta H=0$ and $H$ is a tempered distribution and hence is a polynomial.

Hence we know that $\zeta-g$ is a polynomial as well, however, we~have that $\zeta-g=0$ at $r=0$, hence it is also a polynomial in $r$ only. However, $g\to 0$ for $r\to \infty$ and hence if $\zeta-g$ does not vanish, then $\zeta$ would have to grow at least linearly in $r$. For all $a$ and large enough $R$ we would then have
\begin{align*}
1\lesssim \frac{1}{R}\bigl|\zeta(R,a)-\zeta(\sfrac{R}{2},a)\bigr|\lesssim \frac{1}{R}\int_{\sfrac{R}{2}}^R|\de_r\zeta(s,a)|\ds\lesssim \biggl(\int_{\sfrac{R}{2}}^R\frac{1}{s}|\de_r\zeta(s,a)|^2\ds\biggr)^{\sfrac{1}{2}}.
\end{align*}
By taking a square root and using Fubini's theorem, we~obtain that $\frac{1}{\sqrt{r}}\nabla\zeta\notin L^2$, which is a contradiction.
\end{proof}

\begin{lemma}\label{rot pre}
The Euler equation \eqref{Euler1} holds if \eqref{Euler2} and \eqref{Euler3} hold and the circulations $\gamma_i=\int_{\de B_i}\tau\cdot u\dx$ are conserved in time.
\end{lemma}
In the two-dimensional setting this statement is well-known, see \eg \cite{filho2007vortex}.

\begin{proof}
Indeed the vorticity equation always holds and therefore we have
\[
\curl(\de_t u+(u\cdot \nabla)u)=0.
\]
However not every curl-free field has to be a gradient in $\mathcal{F}$, as~the domain has holes. Using the usual path integral construction, it is easy to see that it is a gradient if
\begin{align*}
\int_\Gamma (\de_tu+(u\cdot \nabla)u)\cdot \tau_\Gamma\dx\eq0
\end{align*}
along every closed, non-self-intersecting path $\Gamma\subset \mathcal{F}$ with normalized tangent $\tau_\Gamma$, which has winding number $1$ with respect to exactly one $B_i$ and $0$ with respect to all others. If~$\Phi_t$ is the flow induced by $u$, then by direct calculation one sees that
\begin{align*}
\de_s\int_{\Phi_s(\Gamma)}u\cdot\tau_{\Phi_s(\Gamma)}\dx\big|_{s=0}\eq\int_\Gamma (\de_tu+(u\cdot \nabla)u))\cdot \tau_\Gamma\dx.
\end{align*}
On the other hand we see that for such $\Gamma$ it holds
\begin{align*}
\int_\Gamma u\cdot\tau_\Gamma\dx\eq\int_{\de B_i} u\cdot\tau\dx,
\end{align*}
which is constant by assumption, and hence we see that there is a $p$ such that \hbox{$\de_t u+(u\cdot\nabla)u=-\nabla p$}.
\end{proof}

\begin{proposition}\label{prop rep u}
The function
\begin{align}\label{u1u2}
u(t)\eq\sum_{i=1}^k\Bigl(\nabla\phi_{i,\dot{q}_i}+\gamma_i\frac{1}{r}\nabla^\perp\psi_i\Bigr)=:u_1+u_2
\end{align}
is a solution to the axisymmetric Euler equations \eqref{Euler1}--\eqref{Euler5}.
\end{proposition}
\begin{proof}
A direct calculation reveals that $u$ is $\curl$-free and fulfills $\div(ru)=0$ and hence $u$ fulfills \eqref{Euler2} and \eqref{Euler3} in $\mathcal{F}$. We~further observe that
\begin{align*}
n\cdot\sum_{i=1}^k\nabla\phi_{i,\dot{q}_i}&\eq u(\dot{q}) \quad\text{on $\ts\bigcup_j\de B_j$},\\
n\cdot\frac{1}{r}\nabla^\perp\psi_i&\eq 0 \quad\text{on $\ts\bigcup_j\de B_j$},\\
\int_{\de B_j}\nabla\phi_{i,\dot{q}_i}\cdot\tau\dx&\eq \int_{\de B_j}\de_\tau\phi_{i,\dot{q}_i}\dx\eq0,\\
\int_{\de B_j}\frac{1}{r}\nabla^\perp\psi_{i}\cdot\tau\dx&\eq \int_{\de B_j}\frac{1}{r}\de_n\psi_{i}\dx\eq\delta_{ij},\\
\frac{1}{r}\de_z\psi_i\eq \de_r\phi_{i,\dot{q}_i}&\eq 0\quad\text{on $\de \mathbb{H}$},
\end{align*}
where $\tau=n^\perp$.

Hence this $u$ has the prescribed circulations and boundary velocities, which shows the statement by Lemma \ref{rot pre}.
\end{proof}

We shall refer to both the $\phi_{i,t^*}$ and their sum as potentials of $u$. We~shall refer to both the $\psi_i$ and their weighted sum as streamfunctions of $u$.

\begin{remark}\label{uunique}
This $u$ is uniquely determined (in $L_R^2$) by $q,\,\dot{q}$ and the $\gamma_i$. Indeed, if~there would be two such different $u$, then their difference would give rise to a nonzero streamfunction with zero circulation (by the same argument as in the proof of Lemma~\ref{ExistencePsi}, Step 3), which is impossible by the uniqueness of the functions $\psi_i$.
\end{remark}

\subsubsection{Representation of $\de_t u$}
We will need to show that the potential and stream function are differentiable in $q$ to be able to represent $\de_tu$. The differentiability of solutions to elliptic equations with respect to changes of the underlying domain is a classical topic and we refer the reader to \cite{sokolowski1992introduction} for further reading.

\skpt
\begin{lemma}\label{SmoothnessInq}
\begin{enumerate}
\item\label{SmoothnessInqa} The function $\phi_{i,\sbullet}$ is smooth as a map from the tangent bundle $TM$ to $H_R^1$ (here differentiability can be understood in both the $L_{\loc}^2$-sense and the pointwise sense).
\item\label{SmoothnessInqb} The derivatives $\de_q \phi_{i,t^*},\de_q^2\phi_{i,t^*}$ lie in $H_R^1\cap C^\infty(\mathcal{F})$, furthermore their $H_{R}^1$-norm is bounded locally uniformly in $q$ and $t^*$.
\item\label{SmoothnessInqc} $\div(r\nabla\de_q \phi_{i,t^*})=0$.
\item\label{SmoothnessInqd}
$|\nabla^m\de_q\phi_{i,t^*}(x)|\lesssim \sfrac{1}{1+|x|^{2+m}}$ for all $m\in\N_{\geq 0}$.
Here the implicit constant is bounded locally uniformly in $q$ and $t^*$.
\end{enumerate}
\end{lemma}

\begin{proof}
We can identify the tangent space at every point with $\R^2$ and $\phi_{i,\sbullet}$ is linear in $t^*$, hence it suffices to show smoothness in $q$ for a fixed $t^*$.

\subsubsection*{Step 1} We first want to apply the implicit function theorem to obtain that a derivative of $\phi_{i,t^*}$ with respect to $q$ exists. Fix some $q^0$. We~use the three-dimensional $\phi^{\R^3}=\phi_{i,t^*}(r,z)$ (in axisymmetrical coordinates) again. We~set
\begin{align*}
\ts
V:= (\dot{H}^1\cap L^6)(\R^3\backslash \bigcup_j B_j^{\R^3}(q^0)),
\end{align*}
and equip this with the standard inner product of $\dot{H}^1$. By~the Sobolev embedding this is a Hilbert space.

In order to fit different configurations of the bodies into one space, we~introduce $C^\infty$ diffeomorphisms $\Xi:\R^3\backslash\bigcup_j B_j(q^0)\to\R^3\backslash \bigcup_j B_j(q^1)$ which map each $\de B_j(q^0)$ to $\de B_j(q^1)$. We~can assume that the family $\Xi$ is smooth in the parameter $q^1$ since the~$B_i$ are. We~may also assume that $\Xi$ is the identity outside a large ball depending on~$q$, but bounded locally uniformly in $q$.
Then $\phi^{\R^3}$ is harmonic on $\R^3\backslash\bigcup_j B_j(q^1)$ with Neumann boundary values $u(t^*,q^1)$ iff the function $\hat{\phi}:=\phi\circ \Xi$ fulfills
\begin{multline}\label{PullbackEq}
\int_{\mathcal{F}(q^0)^{\R^3}}\bscalar{\nabla\hat{\phi}(\text{D}\Xi)^{-1}}{\nabla\eta(\text{D}\Xi)^{-1}}\,|{\det \text{D}\Xi}|\dx\\[-10pt]
\eq-\int_{\cup \de B_j(q^0)^{\R^3}}\frac{\rho_j(q^1)}{\rho_j(q^0)}\,u(t^*,q^1)\eta\dx,
\end{multline}
for all $\eta\in V$, where we have written the minor radius $\rho_j$ as a function of $q$. We~may interpret the difference of the left- and right-hand side as a map $\mathcal{G}:M\times V\to V^*$.

Since $\Xi$ is smooth in $q$ and compactly supported, we~obtain that this map is Fréchet-smooth. Furthermore we have that
\begin{align*}D_V\mathcal{G}(q^0,\hat{\phi})\cdot\delta\phi\eq\int_{\mathcal{F}(q^0)^{\R^3}} \scalar{\nabla\delta\phi}{\nabla\cdot}\dx.
\end{align*}
This is an isomorphism by the Riesz representation theorem. Hence we see that $\hat{\phi}$ is Fréchet-smooth by the implicit function theorem.
This implies that a function $\de_q\phi^{\R^3}$ exists in $V$ by the smoothness of $\Xi$ in $q$. Similarly, higher derivatives must exist. This shows \eqref{SmoothnessInqa}.

\subsubsection*{Step 2} Clearly, $\de_q \phi^{\R^3}$ must be harmonic and hence smooth away from the boundary. To see smoothness up to the boundary we differentiate \eqref{PullbackEq} with respect to $q$ at $q^0$ and obtain that
\begin{multline}
\int_{\mathcal{F}(q^0)^{\R^3}}\bscalar{\nabla\de_q\hat{\phi}(\text{D}\Xi)^{-1}}{\nabla\eta(\text{D}\Xi)^{-1}}\,|{\det \mathrm{D}\Xi}|\\
+\bscalar{\nabla\hat{\phi}\de_q((\text{D}\Xi)^{-1})}{\nabla\eta(\text{D}\Xi)^{-1}}\,|{\det \text{D}\Xi}|
+ \bscalar{\nabla\hat{\phi}(\text{D}\Xi)^{-1}}{\de_q\left(\nabla\eta(\text{D}\Xi)^{-1}\,|{\det \text{D}\Xi}|\right)}\dx\label{EquationDeri}\\
\eq-\int_{\cup \de B_j(q^0)^{\R^3} }\Bscalar{\de_q\Bigl(\frac{\rho_j(q)}{\rho_j(q^0)}u(t^*,q^0)\Bigr)}{\eta}\dx,
\end{multline}
for all $\eta\in V$, the differentiation of this equation is justified by the differentiability of~$\hat{\phi}$ and by $\Xi$ being compactly supported and smooth in $q$.

This is an elliptic equation for $\de_q\hat{\phi}$ with Neumann boundary conditions and a smooth and compactly supported term given by the second and third summand on the left hand side. Hence $\de_q\hat{\phi}$ is smooth up the boundary. The same argument can be used to show regularity of higher derivatives in $q$. Again this also shows that $\de_q\phi^{\R^3}$ (and hence also $\phi_{i,t^*}$) is smooth up to the boundary and the same is true for higher derivatives in $q$.

By the pointwise smoothness that follows from this, it is obvious that \eqref{SmoothnessInqc} holds.

\subsubsection*{Step 3} To obtain the decay estimate for the derivative, we~note that it is enough to show these estimates for $\de_q\phi^{\R^3}$. Clearly, it holds that $\Delta\de_q\phi^{\R^3}=0$.
We again employ Lemma \ref{decay1}, by \eg going back to $\hat{\phi}$ and using Equation \eqref{EquationDeri}, it is easy to see that the boundary values are locally uniformly controlled by $q$ and $t^*$. To see that the integral over the Neumann boundary values of $\de_q\phi^{\R^3}$ is $0$, we~introduce some compact~$B_j'$ with smooth boundary, in which $B_j^{\R^3}$ is compactly contained and which intersects no other $B_{j'}^{\R^3}$. Then we rewrite them as
\begin{align*}
\int_{\de B_j^{\R^3}}\de_n\de_q\phi^{\R^3}\dx\eq -\int_{\de B_j'}\de_n\de_q\phi^{\R^3}\dx\eq-\de_q\int_{\de B_j'}\de_n\phi^{\R^3}\dx\eq0.
\end{align*}
Here pulling out the derivative is justified by the regularity of $\de_q\phi^{\R^3}$.
In particular, the decay estimate also implies that the derivative is in $H^1$.
\end{proof}

\begin{lemma}
The functions $f_i$ are smooth in $R_i$, in particular, $E_{q_i}$ is smooth with respect to $q$.
\end{lemma}
\begin{proof}This can be shown as in the previous lemma by using a similar smooth family of diffeomorphisms.\end{proof}

\skpt
\begin{lemma}\label{PsiDeri}
\begin{enumerate}
\item\label{PsiDeria} The derivative of $\psi_i$ with respect to $q$ exists and is smooth up to the boundary (here the derivative can \eg be taken as a classical pointwise derivative or in the $L_{\loc}^2$ sense).
\item\label{PsiDerib} We have that $\frac{1}{r}\nabla\de_q\psi_i,\,\frac{1}{r}\nabla\de_q^2\psi_i\in L_R^2\cap C^\infty$. Furthermore the $L^2$-norm of these derivatives is bounded locally uniformly in $q$.
\item\label{PsiDeric} It holds that
\[
\left|\de_q\psi_i(x)\right|\lesssim \frac{1}{1+|x|}\qquad\text{and}\qquad
\Bigl|\nabla^m\de_q\frac{1}{r}\nabla^\perp\psi_i(x)\Bigr|\lesssim \frac{1}{1+|x|^{2+m}}\quad \forall m\in\N_{\geq0}.
\]
In the second estimate the implicit constant is locally uniformly bounded in $q$.
\item\label{PsiDerid} The values $C_{ij}$ are differentiable with respect to $q$.
\end{enumerate}
\end{lemma}

\skpt
\begin{proof}

\eqref{PsiDeria} and \eqref{PsiDerid} The argument uses a similar technique as the existence proof for Lem\-ma~\ref{ExistencePsi}. First, we~again consider the three-dimensional vector potentials $\Psi_i$ as in said proof, for fixed boundary values $(\tilde{C}_{ij})$, they have arbitrarily many derivatives in~$q$, which are smooth up to the boundary by the same argument as in the previous proof and the derivatives are in $\dot{H}^1\cap L^6$. This also shows that for fixed $(\tilde{C}_{ij})$ there is a (smooth) derivative of $u_{2,i}$ and $\psi_{i}(\tilde{C}_{ij})$.

It remains to argue that the values $C_{ij}$ are differentiable. To see this note that the linear map from the $(\tilde{C}_{ij})$ to the integrals $\int_{\de B_l}\frac{1}{r}\de_n\psi_i(\tilde{C}_{ij})\dx$, which was used to show existence of the values $C_{ij}$, is differentiable in $q$ as well. Indeed we may again introduce some compact $B_l'$, which compactly contains $B_l$ and intersects no other $B_{j}$. Then we have
\begin{align*}
\int_{\de B_l'}\frac{1}{r}\de_n\de_q\psi_i(\tilde{C}_{ij})\dx\eq\de_q\int_{\de B_l'}\frac{1}{r}\de_n\psi_i(\tilde{C}_{ij})\dx\eq\de_q \int_{\de B_l}\frac{1}{r}\de_n\psi_i(\tilde{C}_{ij})\dx,
\end{align*}
which shows differentiability.
As the $\dot{H}^1$-norm of $\Psi_i$ corresponds to the $L^2$-norm of $\frac{1}{r}\nabla^\perp\psi_i(\tilde{C}_{ij})$, we~see the boundedness statement \eqref{PsiDerib}.

The decay of $\de_q\psi_i$ again follows from the fact that the derivative fulfills
\begin{equation*}\div(\frac{1}{r}\nabla\de_q\psi_i)=0\end{equation*} and is smooth up to the boundary by using Lemma \ref{rep psi}.

The decay of $\de_q \frac{1}{r}\nabla^\perp\psi_i$ follows from the fact that the derivative of the three-dimensional stream function is harmonic and smooth up to the boundary as in the previous proof by Lemma \ref{decay1}, and can be controlled locally uniformly in $q$.
\end{proof}

\begin{remark}
Note that if $q$ is $C^2$ in time and $u$ is a solution of the Euler equations \eqref{Euler1}--\eqref{Euler5}, then we must have
\begin{align}\label{eqDu}
\de_t u\eq \sum_{i=1}^k\de_q \phi_{i,\dot{q}}\cdot\dot{q}+\phi_{i,\ddot{q}_i}+\gamma_i\frac{1}{r}\nabla^\perp\de_q\psi_i\cdot \dot{q}.
\end{align}
Indeed this follows from the fact that $u$ is uniquely determined through $q,\dot{q}$ and the circulations $\gamma_i$ (Remark \ref{uunique}). In~particular, $u$ has the regularity required in Condition \ref{strongsol}.

\end{remark}

\begin{lemma}\label{Decayp}
Assume that $u$ solves the Euler equations \eqref{Euler1}--\eqref{Euler5} with pressure~$p$.
Then we have that
\begin{align*}|\nabla p(x)|\lesssim \frac{1}{1+|x|^2},
\end{align*}
and there is a constant $C$ which may be chosen as $0$ such that
\begin{align*}|p(x)-C|\lesssim \frac{1}{1+|x|}.
\end{align*}
The implicit constants in these estimates are bounded locally uniformly in $q,\dot{q},\ddot{q}$.
\end{lemma}
\begin{proof}
By the construction of $u$ in \eqref{u1u2} and the decay estimates in Lem\-mas~\ref{ExistencePhi} and \ref{ExistencePsi} we have that
\begin{align*}
|(u\cdot\nabla)u(x)|\lesssim \frac{1}{1+|x|^5}.
\end{align*}
By Equation \eqref{eqDu} and Lemmas \ref{ExistencePhi}, \ref{SmoothnessInq} and \ref{PsiDeri}, we~see that
\begin{align*}
|\de_tu(x)|\lesssim \frac{1}{1+|x|^2}.
\end{align*}
Hence $|\nabla p(x)|\lesssim \spfrac{1}{1+|x|^2}$ and the estimate is locally uniform in $q,\dot{q},\ddot{q}$, because the estimates for $u$ and its derivatives are.

Since $\int_{\de B_R(0)\cap \mathbb{H}}|\nabla p|\dx\to 0$ for $R\to\infty$ we have
\begin{align*}
\max_{x\in \de B_R(0)} p(x)-\min_{x\in \de B_R(0)} p(x)\to 0.
\end{align*}
Furthermore $\int_1^\infty |\nabla p(x,a)|\dx<\infty$ for all $a$ for which no $B_i$ intersects this line, hence we obtain that $p$ converges to some finite value at infinity, which then gives the decay statement for $p$ by the fundamental theorem of calculus.
\end{proof}

\subsection{Derivation of an ODE for the system}
We reduce the motion of the bodies~$B_i$ to an ODE whose coefficients depend on the functions $\psi_i$ and $\phi_i$. This will also yield existence and uniqueness of solutions to the system.
In two-dimensional bounded domains, a similar calculation can be found \eg in \cite{GlassMunnierSueur2}.

We first introduce some additional terminology.
We set $\phi(t^*)=\sum \phi_{i,t_i^*}$ if $t^*=t_1^*+\dots+ t_k^*$. Furthermore we set $\psi=\sum_i\gamma_i\psi_i$.

\begin{definition}\label{main def}
Let $t^*=t_1^*+\dots+ t_k^*$; $s^*=s_1^*+\dots+ s_k^*$ and $w^*=w_1^*+\dots+w_k^*$ for $t_i^*,s_i^*$ and $w_i^*$ associated with $B_i$. We~define
\begin{align*}
G_i(q,\gamma)\cdot t_i^*&\eq\int_{\de B_i} \frac{1}{2r}\left((\de_n\psi)^2\de_n\phi_{i,t_i^*}\right)\dx,\\
(\mathcal{M}_{ij}(q)t_i^*)\cdot s_j^*&\eq\int_{\mathcal{F}}r\nabla\phi_{i,t_i^*}\nabla \phi_{j,s_j^*}\dx,\\
\langle \Gamma(q),t^*,s^*\rangle\cdot w^*&\eq \frac{1}{2}\sum_{ij}\Bigl( \bigl(\left(\de_q\mathcal{M}_{ij}\cdot s^*\right)t^*\bigr)\cdot w^*+\bigl(\left(\de_q\mathcal{M}_{ij}\cdot t^*\right)s^*\bigr)\cdot w^*,\\[-8pt]
&\hspace*{5cm}-\bigl(\left(\de_q\mathcal{M}_{ij}\cdot w^*\right)s^*\bigr)\cdot t^*\Bigr),\\
(A(q,\gamma)t^*)\cdot s^*&\eq\sum_i\int_{\de B_i}\Bigl(-\de_\tau\phi(s^*)\de_n\phi(t^*)+\de_\tau\phi(t^*)\de_n\phi(s^*)\Bigr)\de_n\psi\dx,
\end{align*}
where in the definition of $\Gamma$, the inner dot product refers to the derivative in that direction.

Furthermore, $\mathcal{M}$ shall be the matrix made up of the blocks $\mathcal{M}_{ij}$ and $E$ shall be the diagonal matrix made up of the blocks $E_{q_i}$. Let $G\in (\R^2)^k\simeq \R^{2k}$ be the vector with the entries $G_1,\dots G_k$.

\end{definition}

\begin{theorem}\label{ODEtheorem}
The system detailed in the Introduction is equivalent to the system of ODEs given by
\begin{equation}\label{MainODE}
E(q)\ddot{q}+\frac{1}{2}\dot{q}(\de_{q}E(q)\cdot\dot{q})+\mathcal{M}(q)\ddot{q}+\langle\Gamma(q),\dot{q},\dot{q}\rangle
\eq G(q,\gamma)+(A(q,\gamma)\dot{q}).
\end{equation}
\end{theorem}
\begin{remark}
The equation can be interpreted as the geodesic equation for the metric given by $\mathcal{M}+E$, with extra terms due to the circulation on the right hand side. The matrix $\mathcal{M}$ describes the ``added inertia'', which encodes the fact that to accelerate one of the bodies, one also has to accelerate the surrounding fluid.

\end{remark}
\begin{proof}
We argued in Remark \ref{uunique} that $u$ is uniquely determined by $q,\dot{q}$, hence it suffices to show that the family of equations in \eqref{2ndDeri} is equivalent to this system. Let~$t_i^*$ be an arbitrary tangent vector associated with $B_i$. We~set $u_i^*=\nabla(\phi_{i,t_i^*})$.

Then by Equation \ref{2ndDeri} it holds that
\begin{align*}
(t_i^*)^TE_{q_i}\ddot{q}_i+\frac{1}{2}\dot{q}_i^T(\de_{q_i}E_{q_i}\cdot\dot{q}_i)t_i^*\eq -\int_{\de B_i} rpu_i^*\cdot n\dx.
\end{align*}
By the equation for $p$, partial integration and the identity $u\nabla u=\frac{1}{2}\nabla|u|^2+u\cdot \curl u$ it follows that
\begin{align*}
\mel(t_i^*)^TE_{q_i}\ddot{q}_i+\frac{1}{2}\dot{q}_i^T(\de_{q_i}E_{q_i}\cdot\dot{q}_i)t_i^*&\eq\int_{\mathcal{F}}r\nabla p\cdot u_i^*\dx\\
&\eq-\int _{\mathcal{F}}r\Bigl(\de_t(u_1+u_2)+\frac{1}{2}\nabla|u_1+u_2|^2\Bigr)\cdot u_i^*\dx.
\end{align*}
It follows from the decay estimates in Lemmas \ref{ExistencePhi}, \ref{ExistencePsi} and \ref{Decayp} that there are no boundary terms from $\infty$ in this partial integration.

We now split this into the different contributions and use the proposition below to obtain the equation in the theorem, tested against $t_i^*$. Since $t_i^*$ was arbitrary this implies the statement.\end{proof} 
\skpt
\begin{proposition}\label{Coefficients}
\begin{enumerate}
\item\label{Coefficientsa}
We have
\begin{align*}
-\frac{1}{2}\int_\mathcal{F} r\nabla|u_2|^2\cdot u_i^*\dx\eq G_i(q,\gamma)\cdot t_i^*.
\end{align*}
\item\label{Coefficientsb}
It holds that
\begin{align*}
-\int_\mathcal{F} r\bigl(\de_tu_2+\nabla(u_1\cdot u_2)\bigr)\cdot u_i^*\dx\eq (A(q,\gamma)\dot{q})\cdot t_i^*.
\end{align*}
\item\label{Coefficientsc}
We have that
\begin{align*}
\int_\mathcal{F} r\Bigl(\de_t u_1+\frac{1}{2}\nabla|u_1|^2 \Bigr)\cdot u_i^*\dx\eq \ddot{q}^T\mathcal{M}(q)t_i^*+\scalar{\Gamma(q),\dot{q}}{\dot{q}}\cdot t_i^*.
\end{align*}
\end{enumerate}
\end{proposition}

\skpt
\begin{proof}

\eqref{Coefficientsa}
Using that both $u_2$ and $u_i^*$ decay like $\sfrac{1}{|x|^2}$ by Lemma \ref{ExistencePhi} and Lemma~\ref{ExistencePsi} we may partially integrate the left-hand side to obtain equality with
\begin{align}
\int_{\de \mathcal{F}}\frac{1}{2}r|u_2|^2\de_n \phi_{i,t^*}\dx.
\end{align}
To see that this equals the definition of $G$ we note that $\de_n \phi_{i,t_i^*}$ vanishes on every boundary except $\de B_i$ and that $|u_2|=\frac{1}{r}|\nabla^\perp\sum_j \gamma_j\psi_j|=\frac{1}{r}|\de_n \sum_j \gamma_j\psi_j|$ since the tangential derivative of the functions $\psi_j$ vanishes.

\eqref{Coefficientsb}
We have that
\begin{align*}
-\int_\mathcal{F} r\nabla(u_1\cdot u_2)u_i^*\dx\eq\int_{\de \mathcal{F}}r(u_1\cdot u_2)(u_i^*\cdot n)\dx,
\end{align*}
(this partial integration is justified by the decay estimates from Lemmas \ref{ExistencePhi} and~\ref{ExistencePsi}) which by the construction of $u$ in \eqref{u1u2} equals
\begin{align*}
\sum_{l}\int_{\de B_i}r\Bigl(\frac{1}{r}\de_n\sum_j\gamma_j\psi_j\Bigr)(\de_\tau \phi_{l,\dot{q}_l})(\de_n \phi_{i,t_i^*})\dx,
\end{align*}
because $u_2$ has no normal component on the boundary.

It holds that $\de_tu_2=\frac{1}{r}\nabla^\perp\de_t\psi$. We~have that $\div(\frac{1}{r}\nabla\de_t\psi)=0$ and $\de_t\psi$ has the boundary values
\begin{align*}
\de_t\psi\eq\sum_j\gamma_j\de_{q} C_{jl}\cdot\dot{q}-(u_1\cdot n)\de_n\psi\quad\text{on $\de B_l$}
\end{align*}
as one can see by differentiating the identity $C_{jl}(q)=\psi_j(x_q)(q)$ where $x_q$ is some fixed point on $\de B_l$ whose derivative in $t$ equals $u_1\cdot n$.
Then a partial integration, which can again be justified by the decay estimates in Lemmas \ref{ExistencePhi} and \ref{PsiDeri}, reveals that
\begin{align*}
-\int_{\mathcal{F}}ru_i^*\cdot\de_tu_2\dx\eq\sum_l\sum_j\int_{\de B_l} \de_\tau\phi_{i,t_i^*}\left(\gamma_j\de_{q} C_{jl}\cdot\dot{q}-(u_1\cdot n)\de_n\psi\right)\dx.
\end{align*}
The first summand vanishes because $\de_q C_{jl}$ is a constant on each $\de B_l$ and this proves~\eqref{Coefficientsb}.

\eqref{Coefficientsc}
We introduce an energy functional for the potential part of the fluid velocity:
\begin{align*}
\mathcal{E}_{u_1}:=\frac{1}{2}\int_{\mathcal{F}} r|u_1|^2\dx.
\end{align*}
Following the approach in \cite{Munnier5}, we~will show that
\begin{align}\label{Claim1}
(\de_t\de_{\dot{q}}-\de_{q})\mathcal{E}_{u_1}\cdot t_i^*\,=\int_{\mathcal{F}}u_i^*\cdot(\de_tu_1+\frac{1}{2}\nabla|u_1|^2)\dx.
\end{align}
To prove this claim we shall use the following lemma which can be proved exactly as in \cite[Lem.\,5.1]{Munnier5}: \begin{lemma}\label{deriForm}
For $\eta\in \dot{H}_R^1$ let
\begin{align*}&\Lambda(\eta)\eq\int_{\mathcal{F}}r\scalar{\nabla\phi(\dot{q})}{\nabla \eta} \dx.
\end{align*}
Then it holds that
\begin{align*}
(\de_t\de_{\dot{q}}-\de_q)(\Lambda)\eq0.
\end{align*}\end{lemma}
We now note that
$\mathcal{E}_{u_1}\eq\frac{1}{2}\Lambda(\phi(\dot{q}))$ and that
\begin{align*}
\de_{\dot{q}}\mathcal{E}_{u_1}\cdot t_i^*\eq\frac{1}{2}\bigl((\de_{\dot{q}}\Lambda)(\phi(\dot{q}))\cdot t_i^*+\Lambda(\phi(t_i^*))\bigr).
\end{align*}
Because $\phi(t_i^*)$ also equals $\de_{\dot{q}}\phi(\dot{q})\cdot t_i^*$, we~see that
\begin{align}\label{eq 5}
\de_{\dot{q}}\mathcal{E}_{u_1}\cdot t_i^*\eq(\de_{\dot{q}}\Lambda)(\phi(\dot{q}))\cdot t_i^*.
\end{align}
Hence we obtain that
\begin{multline*}
(\de_t\de_{\dot{q}}-\de_{q})\mathcal{E}_{u_1}\cdot t_i^*\\
\eq\left(\de_t\de_{\dot{q}}\cdot t_i^*-\de_q\cdot t_i^*\right)\Lambda(\phi)
+(\de_{\dot{q}}\Lambda)(\de_t\phi(\dot{q}))\cdot t_i^*
+\frac{1}{2}(\de_q\Lambda)(\phi(\dot{q}))\cdot t_i^*
-\frac{1}{2}\Lambda(\phi^\dagger),
\end{multline*}
where $\phi^\dagger=\de_q\phi(\dot{q})\cdot t_i^*$ and we made use of \eqref{eq 5}.

By Lemma \ref{deriForm}, the first term is $0$. By~definition, the second term equals
\begin{align*}
\de_{\dot{q}}\Lambda(\de_t\phi(\dot{q}))\cdot t_i^*\eq\int_{\mathcal{F}}r\de_tu_1\nabla \phi(t_i^*)\dx.
\end{align*}
By the Reynolds transport theorem (whose usage is justified by Lemma \ref{SmoothnessInq}\,\eqref{SmoothnessInqb}), we~have that twice the third term equals
\begin{align*}
(\de_q\Lambda\cdot t_i^*)(\phi)\eq-\int_{\de B_i}r|u_1|^2u_i^*\cdot n\dx+\Lambda(\phi^\dagger).
\end{align*}
This yields the claim \eqref{Claim1} after another partial integration.

Now we use that $\mathcal{E}_{u_1}=\frac{1}{2}\mathcal{M}(q)\dot{q}\cdot\dot{q}$, which follows directly from the definition of $\mathcal{M}$. Then we may compute the Euler-Lagrange equation of this as
\begin{align*}
(\de_t\de_{\dot{q}}-\de_{q})\mathcal{E}_{u_1}\cdot t_i^*\eq \mathcal{M}(q)\ddot{q}\cdot t_i^*+((\de_q\mathcal{M}(q)\cdot \dot{q})\dot{q})\cdot t_i^*-\frac{1}{2}((\de_q\mathcal{M}(q)\cdot t_i^*)\dot{q})\cdot\dot{q}.
\end{align*}
The last two summands equal the Christoffel symbol $\Gamma$ as one can directly see by writing them out in components.
\end{proof}
\begin{thebibliography}{BGR23}
\bibitem[11]{evans2022partial} L. C. Evans – Partial differential equations, second ed., Graduate Studies in Math., vol. 19, American Mathematical Society, Providence, RI, 2010.
\bibitem[14]{GlassMunnierSueur2} O. Glass, C. Lacave, A. Munnier \& F. Sueur – “Dynamics of rigid bodies in a two dimensional incompressible perfect fluid”, J. Differential Equations 267 (2019), no. 6, p. 3561–3577.
\bibitem[1]{Amrouche} C. Amrouche, V. Girault \& J. Giroire – “Dirichlet and Neumann exterior problems for the n-dimensional Laplace operator: An approach in weighted Sobolev spaces”, J. Math. Pures Appl. (9) 76 (1997), no. 1, p. 55–81.
\bibitem[19]{grisvard2011elliptic} P. Grisvard – Elliptic problems in nonsmooth domains, Classics in Applied Math., vol. 69, Society for Industrial and Applied Mathematics (SIAM), Philadelphia, PA, 2011.
\bibitem[13]{GallaySverak} T. Gallay \& V. Šverák – “Remarks on the Cauchy problem for the axisymmetric Navier-Stokes equations”, Confluentes Math. 7 (2015), no. 2, p. 67–92.
\bibitem[12]{filho2007vortex} M. L. Filho – “Vortex dynamics in a two-dimensional domain with holes and the small obstacle limit”, SIAM J. Math. Anal. 39 (2007), no. 2, p. 422–436.
\bibitem[35]{sokolowski1992introduction} J. Sokołowski \& J.-P. Zolésio – Introduction to shape optimization, Springer Series in Computational Math., vol. 16, Springer-Verlag, Berlin, 1992.
\bibitem[33]{Munnier5} A. Munnier – “Locomotion of deformable bodies in an ideal fluid: Newtonian versus Lagrangian formalisms”, J. Nonlinear Sci. 19 (2009), no. 6, p. 665–715.
\end{thebibliography}
\end{document}
